\documentclass[11pt]{article}
\usepackage[margin=1in]{geometry}
\usepackage{amsmath,amssymb,amsthm,booktabs}
\newtheorem{theorem}{Theorem}
\begin{document}
\title{An explicit optimal three-point flag SDP for the first plateau}
\author{}\date{}\maketitle
Put
\[
s=\sqrt7,\qquad t=\frac{s-1}{2},\qquad
B=2s-5,\qquad V=\frac{2s-4}{3},\qquad a=3V=2s-4.
\]
We give a single rank-one matrix attaining the exact first plateau.
It proves the sharp graph inequality directly, without an edge-elimination
or fixed-point reduction in the upper-bound argument.

\paragraph{The five patterns and four flags.}
Consider an ordered graph with the interval property: for $x<y<z$,
an edge $xz$ implies both $xy$ and $yz$.
Its three-point patterns are $0$ (empty), $L$ (only the lower pair),
$H$ (only the upper pair), $P$ (both consecutive pairs), and $T$ (triangle).
For a single root, order the two-point flags as
\[
(F_1,F_2,F_3,F_4)
=(\text{lower nonneighbor},\text{lower neighbor},
  \text{upper neighbor},\text{upper nonneighbor}).
\]
Let $M(G)$ be the flag-product matrix for a three-point pattern $G$:
choose the root uniformly among its three vertices, and choose the two
remaining vertices as the two disjoint one-vertex extensions, in random
order. This is the standard family $(t_{\rm root},f)=(1,2)$ at sample
size three.

\begin{theorem}[An explicit optimal matrix]
For the graph objective $g(G)=(1+\mathbf1_H(G)-\mathbf1_T(G))/3$,
the three-point flag SDP has optimal upper bound $V$.
A single optimal PSD matrix is
\[
\boxed{
Q=Vvv^\top,
\qquad v=(-t,1,1,0)^\top,
}
\]
or explicitly
\[
Q=\begin{pmatrix}
2s-5&s-3&s-3&0\\
s-3&V&V&0\\
s-3&V&V&0\\
0&0&0&0
\end{pmatrix}.
\]
\end{theorem}
\begin{proof}[Certificate]
Since $V>0$, the displayed rank-one factorization proves $Q\succeq0$.
Direct averaging over the root gives
\[
\begin{array}{c|ccccc}
G&0&L&H&P&T\\\hline
v^\top M(G)v&t^2/3&t^2/3&-2t/3&(1-t)/3&1
\end{array}.
\]
Using $2t^2+2t-3=0$ and $V=2/(3+2t)$, the complete list of
certificate values is
\[
\boxed{
g(G)+\langle Q,M(G)\rangle=
\begin{cases}
V,&G\in\{0,L,H,T\},\\
5V/6,&G=P.
\end{cases}}
\]
Conditional independence of the two extensions implies that the
expectation of the flag term is nonnegative. Consequently $\mathbb E g\le V$.
\end{proof}

\paragraph{The underlying sum-of-squares identity.}
For an iid root $X$, define
\[
d(X)=\mathbb P(Y\sim X\mid X),\qquad
\ell(X)=\mathbb P(Y<X,\ Y\not\sim X\mid X).
\]
The flag vector satisfies $v^\top h(X)=d(X)-t\ell(X)$.
Thus the certificate is equivalently the exact identity
\[
\boxed{
B-(p_H-p_T)
=a\,\mathbb E[(d(X)-t\ell(X))^2]+\frac a6\,p_P.
}
\]
This is a flag sum of squares plus a nonnegative pattern density.
For finite weighted graphs, repeated copies of a vertex are adjacent;
the same identity holds, since the two adjacent flags have equal
coefficients and their relative order does not affect the score.

\paragraph{How to find the matrix analytically.}
The extremal construction consists of ordered disjoint cliques with masses,
from top downward,
\[
q,\ q(1-q),\ q(1-q)^2,\ldots,
\qquad q=\frac{4-\sqrt7}{3}.
\]
At a root in the $k$th clique, its neighbor mass is $q(1-q)^k$, and
its lower nonneighbor mass is $(1-q)^{k+1}$. Their ratio is
$q/(1-q)=t$. This suggests the vanishing flag combination $d-t\ell$
and hence the rank-one direction $v$.

Alternatively, start with $Q_J=a vv^\top$ for the objective $J=p_H-p_T$.
Demand equal certificate values on $0,H,T$, all of which occur in the
extremal construction:
\[
B=\frac{at^2}{3}=1-\frac{2at}{3}=a-1.
\]
These equations give $2t^2+2t-3=0$, its positive root
$t=(\sqrt7-1)/2$, and $a=2\sqrt7-4$.
The remaining pattern $L$ has the same value as $0$, while $P$ has
positive slack $a/6$. Dividing the matrix by three gives $Q$ above.
This constructs the optimal solution without numerical optimization.

\paragraph{Sharpness.}
For the geometric clique construction, $p_P=0$ and $d=t\ell$ almost
surely, so the identity gives $p_H-p_T=B$ and $\mathbb E g=V$.
This is also approached by finite weighted graphs. If the first $N$
cliques retain the displayed masses and the remaining mass $(1-q)^N$
is put in one bottom clique, the defect is
\[
B-(p_H-p_T)=a(1-q)^{3N}\longrightarrow0.
\]
Hence every valid flag SDP upper bound for this graph objective is at
least $V$, and the displayed matrix is optimal.

\paragraph{Application to the additive event.}
Let $1<c\le\sqrt2$, choose $c\le r\le2/c$, and join positive values
$x<y$ when $y<rx$. This graph has the interval property. For a sorted
triple $x<y<z$, let
\[
\sigma=\mathbf1\{x+y<cz\}+\mathbf1\{x+z<cy\}.
\]
A middle win implies $z<cy\le ry$ and
$x<(c-1)y\le y/r$, because $r(c-1)\le1$.
It therefore requires type $H$. A triangle gives $x>z/r$ and
$x+y>2z/r\ge cz$, so it has no wins. It follows that
\[
\sigma\le1+\mathbf1_H-\mathbf1_T.
\]
The actual additive objective is $\sigma/3$; the same matrix $Q$
therefore certifies $S_2(c)\le V$. Retaining the additive labels refines
the five graph patterns but does not change the matrix or the proof.
Positive atomless perturbations extend the upper bound to arbitrary
nonnegative laws, preserving strict winning events in the limit inferior.

For the matching lower bound, put masses $q(1-q)^k$ at $\varepsilon^k$,
where $0<\varepsilon<\min\{c/2,c-1,1/c,1/r\}$.
Within the range $1<c\le\sqrt2$, three equal-scale draws never win;
a largest-scale target wins against two lower-scale draws, and either
of two largest-scale targets wins when the third draw is smaller.
The success probability satisfies
\[
S=(1-q)^3S+q(1-q)^2+2q^2(1-q),
\]
whose solution is $S=V$. Thus the colored/additive three-point SDP itself
is optimal for every $1<c\le\sqrt2$, and
\[
S_2(c)=\frac{2\sqrt7-4}{3}.
\]
The admissible interval $[c,2/c]$ for the ratio cutoff closes at $\sqrt2$.

\paragraph{Exact verification.}
Run
\begin{verbatim}
python3 graph_experiment/verify_first_plateau_flag.py
\end{verbatim}
The standard-library checker enumerates all five interval-graph patterns,
checks the matrix entries and every score, and verifies the sharpness
recurrences in the exact field generated by the root $a\in(1,2)$ of
$a^2+8a-12=0$. The matrix entries lie in $\mathbb Q(\sqrt7)$; positive
semidefiniteness follows from a positive algebraic scalar times a rank-one
outer product. Its JSON output is \texttt{first\_plateau\_flag.json}.
\end{document}
